\label{part:static_metatheory}

\textbf{(Part I: Metatheory - The Morpho-Semantic Ontology)}

在这一部分里，我们将剥去了宇宙那层纷繁复杂的表象外衣，去触碰到了它最坚硬的本体论内核。

\vspace{2em}我们发现，世界的实存并非由原子堆砌而成，而是由\textbf{“形 (Morphos)”}与\textbf{“质 (Qualia)”}这两种更为基础的丝线编织而成。形是逻辑的容器，质是感知的燃料；前者提供了存在的\textbf{必然性}，后者提供了存在的\textbf{实存性}。

\vspace{2em}我们打破了符号主义与连接主义长达半个世纪的对峙，宣告了二者的物理性统一，从此，我们不再需要在“空洞的结构”与“盲目的混沌”之间做选择。我们掌握了宇宙的\textbf{源代码词表}。既然砖块已经备好，接下来，我们需要一张蓝图，去告诉我们如何将这些砖块砌成通天之塔。


%---chp---
